% Преамбула pandoc: подключается через --include-in-header
% Язык (russian) и кодировка (UTF-8) задаются в YAML-метаданных — pandoc сам подключает babel
% Не дублируем fontenc/inputenc/babel — они конфликтуют с шаблоном pandoc

% Абзацный отступ
\usepackage{indentfirst}
\setlength{\parindent}{0.75cm}
\setlength{\parskip}{0pt plus 1pt}

% Поля (ГОСТ для методических пособий)
\usepackage{geometry}
\geometry{top=2cm, bottom=2cm, left=3cm, right=1.5cm}

% Межстрочный интервал (полуторный)
\usepackage{setspace}
\onehalfspacing

% Компактные заголовки и абзацы. scrreprt по умолчанию оставляет слишком много воздуха.
\KOMAoptions{parskip=false}
\RedeclareSectionCommand[beforeskip=-1.0\baselineskip,afterskip=0.45\baselineskip]{chapter}
\RedeclareSectionCommand[beforeskip=-0.7\baselineskip,afterskip=0.25\baselineskip]{section}
\RedeclareSectionCommand[beforeskip=-0.5\baselineskip,afterskip=0.15\baselineskip]{subsection}
\RedeclareSectionCommand[beforeskip=-0.4\baselineskip,afterskip=0.10\baselineskip]{subsubsection}

% Sans-serif шрифт для заголовков (scrreprt использует \sffamily по умолчанию)
\setsansfont{DejaVu Sans}

% Оформление кода: fvextra корректно сохраняет UTF-8/кириллицу в отличие от listings.
\usepackage{xcolor}
\definecolor{codebg}{gray}{0.94}
\definecolor{codeframe}{gray}{0.5}
\definecolor{bashbg}{RGB}{235,240,248}
\definecolor{bashframe}{RGB}{80,120,180}
\usepackage{fvextra}

% Выделенные блоки: предупреждения и дополнительная информация
\usepackage[most]{tcolorbox}
\definecolor{attentionbg}{RGB}{255,250,225}
\definecolor{attentionframe}{RGB}{214,158,0}
\definecolor{additionalbg}{RGB}{238,248,238}
\definecolor{additionalframe}{RGB}{92,145,92}
\newcommand{\attentionicon}{{\color{attentionframe}\sffamily\bfseries ⚠}}
\newtcolorbox{attentionbox}{
  enhanced,
  breakable,
  colback=attentionbg,
  colframe=attentionframe,
  colbacktitle=attentionbg,
  coltitle=black,
  boxrule=0.7pt,
  arc=2pt,
  left=8pt,
  right=8pt,
  top=5pt,
  bottom=6pt,
  before skip=7pt,
  after skip=7pt,
  title={\attentionicon\quad Внимание},
  fonttitle=\bfseries\sffamily
}
\newtcolorbox{additionalbox}{
  enhanced,
  breakable,
  colback=additionalbg,
  colframe=additionalframe,
  colbacktitle=additionalbg,
  coltitle=black,
  boxrule=0.6pt,
  arc=2pt,
  left=8pt,
  right=8pt,
  top=5pt,
  bottom=6pt,
  before skip=7pt,
  after skip=7pt,
  title={Дополнительная информация},
  fonttitle=\bfseries\sffamily
}

% Подписи к таблицам и рисункам
\usepackage{caption}
\captionsetup[table]{
  labelsep=endash,
  justification=raggedright,
  singlelinecheck=false,
  skip=0pt
}
\captionsetup[figure]{
  labelsep=endash,
  justification=centering,
  skip=3pt
}

% Графика
\usepackage{graphicx}
\usepackage{float}
\graphicspath{{./pictures/}}
\makeatletter
\def\maxwidth{\ifdim\Gin@nat@width>\linewidth\linewidth\else\Gin@nat@width\fi}
\def\maxheight{\ifdim\Gin@nat@height>0.48\textheight 0.48\textheight\else\Gin@nat@height\fi}
\setkeys{Gin}{width=\maxwidth,height=\maxheight,keepaspectratio}
\def\fps@figure{H}
\makeatother
\setlength{\textfloatsep}{6pt plus 2pt minus 2pt}
\setlength{\floatsep}{6pt plus 2pt minus 2pt}
\setlength{\intextsep}{6pt plus 2pt minus 2pt}

% Гиперссылки
\usepackage{hyperref}
\hypersetup{
  colorlinks=true,
  linkcolor=black,
  urlcolor=blue,
  citecolor=black
}

% Поддержка кириллицы в закладках PDF
\usepackage{bookmark}

% Таблицы
\usepackage{booktabs}
\usepackage{longtable}

% Настройка оглавления: меньший шрифт, минимальный отступ сверху
\usepackage{tocloft}
\setlength{\cftbeforetoctitleskip}{0em}
\setlength{\cftaftertoctitleskip}{1em}
\renewcommand{\cfttoctitlefont}{\normalfont\large\bfseries}
\renewcommand{\cftchapfont}{\small}
\renewcommand{\cftsecfont}{\footnotesize}
\renewcommand{\cftchappagefont}{\small}
\renewcommand{\cftsecpagefont}{\footnotesize}
\setlength{\cftbeforechapskip}{3pt}
\setlength{\cftbeforesecskip}{1pt}
\AtBeginDocument{\setcounter{tocdepth}{0}}
